% !TeX root = ../main.tex
\section{Introduction}
\label{sec:introduction}

%% FLAG (see manuscript/CLAUDE.md §0.7): the wording below is deliberately neutral about how many
%% stations the study covers, because the extension to three buoys with contrasting wave climates
%% is planned but not yet run. Revisit every sentence marked SITE-COUNT once that is settled.
%% FLAG: citation commands here use \citep/\citet (natbib author-year, correct for elsarticle's
%% authoryear option). sections/02_methods.tex currently uses plain \cite{} throughout — the two
%% should be made consistent before submission.

Short-term forecasts of the sea state underpin a wide range of offshore and coastal activity,
from the planning of installation and offloading windows to the assessment of extreme loads and
lifetime fatigue on marine structures \citep{orimoloye2019bimodal}. Operationally, such forecasts
are produced by numerical wave models (NWM), which integrate the wave energy balance equation
forward in time given wind forcing, boundary conditions and bathymetry. They are conventionally
expressed, distributed and verified in terms of a small number of integrated (bulk) parameters, chiefly the significant wave height $H_s$ and a characteristic period. This convention is deeply
embedded in operational practice. Routine intercomparison of wave forecasts has long been
established under the World Meteorological Organization (WMO), which now compares eighteen regional and global forecast systems verifying them against in situ buoy observations
\citep{ecmwf2026lcwfvproject}. The fields exchanged for that verification are six integrated
parameters: significant wave height, peak and mean zero-crossing periods, mean wave direction and the two components of the 10-m wind.

Nevertheless, the bulk parameters are summaries of a richer quantity, and they are not a
sufficient statistic for the purposes to which a forecast is ultimately put. The physically
complete description of a sea state is the directional wave spectrum, from which every
bulk parameter can in turn be derived, whereas the reverse does not hold \citep{holthuijsen2007}.
Because the response of a coastal defence or a marine structure depends on \emph{where in
frequency} the energy is situated, and not only on how much energy is present, two sea states
carrying identical total energy --- and therefore identical $H_s$ --- need not behave alike \citet{orimoloye2019bimodal}. The point is sharpest in mixed seas, with double or multi-peaked spectra which no single height-period pair can represent \citep{portilla2009,violantecarvalho2002}.

The directional spectrum, although complete, is poorly suited to routine comparison against
observations. Operational buoys resolve only its first five directional moments, and existing estimators cannot accommodate more than two directional peaks and may introduce spurious energy \citep{gorman2018}. The frequency spectrum
$E(f)$ carries no such ambiguity. It is estimated from the heave record alone, as $H_s$ is, and
three of the four wave parameters exchanged for the WMO intercomparison ($H_s$, $T_p$ and $T_z$)
are derived from it alone. Its absence from routine intercomparison reflects the sampling
variability of individual spectral estimates and the absence of a standard spectral skill score
in operational practice, rather than a limitation of the observation; where verification has
been carried out at the level of spectral components, it has revealed model deficiencies that
bulk statistics conceal \citep{hanson2009pacific}. $E(f)$ cannot separate wave systems that share a
frequency band while travelling in different directions, but it resolves swell and wind sea
whenever they are separated in frequency, and it retains the information identified
above as lost in the bulk parameters. It is therefore the most complete description of the sea
state that a buoy observes without additional assumptions.

Independently of how the forecast is represented, the means of producing it has been changing. In
meteorology, data-driven models are now evaluated
head-to-head against operational numerical weather prediction systems under shared ground
truth and shared metrics \citep{rasp2024weatherbench2}, and a first assessment in an
operational-like setting reports accuracy comparable to that of a leading operational system for
both global metrics and extreme events \citep{benbouallegue2024rise}. To our knowledge, comparable gains have been
slower to materialise in ocean forecasting [INSERT: reference is likely needed]. The motivation for pursuing them is nonetheless
substantial, and is chiefly one of cost: a NWM requires a full forcing chain and
considerable computation, whereas a trained model maps recently observed conditions directly to a
forecast, and has been shown to reproduce a numerical model's output at a small fraction of its
computational expense \citep{james2018wave}, prompting the question of whether such models can act
as physical-model surrogates at buoy scale \citep{minuzzi2023lstm}. By contrast with a NWM, however, a data-driven model learns the evolution of the sea state (swell dispersion,
wind-sea growth and decay) implicitly from historical records rather than from an explicit
physical parameterisation, and so obtains its speed at the expense of interpretability and of the
physical guarantees a NWM provides by construction.

Within ocean forecasting, data-driven work has concentrated overwhelmingly on the scalar bulk
parameters rather than on the spectrum itself. From
early artificial neural networks (ANNs) forecasting wave height and period
\citep{makarynskyy2004}, through decomposition-based recurrent approaches to wave-height
forecasting \citep{zhou2021emdlstm, ding2023eofeemd}, to recent transformer architectures
\citep{kim2026metoformer}. Work that does engage with the spectrum generally addresses a different
problem. One line estimates a spectrum, or its parametric descriptors, from bulk parameters or
from measured platform response, rather than forecasting it forward in time
\citep{naithani2005ann, namekar2006ann, sakhare2009svr, nielsen2023shipbuoy, kwon2026transformer}.
A second reconstructs or predicts spectra over a spatial domain from atmospheric forcing
\citep{song2023jpo, gao2025jgr, liu2025cnnxlstm, cao2025bohai}. Forecasting the spectrum forward
in time from a station's own recent spectral history has received considerably less attention.

Before proposing such a model, a further caution applies. In a comparison across sixteen buoys,
\citet{jiang2024comment} found that autoregression, gradient boosting and several deep
architectures performed almost identically for univariate significant-wave-height forecasting,
and concluded that the complex models mostly recover the linear autocorrelation a simple baseline
already exploits. We consider this critique sound and it shapes the present
study in two ways. First, every forecast is evaluated against a linear and a persistence baseline,
rather than assuming the more complex learned model has the advantage. Second, it leaves open whether the same conclusion holds when the target is the spectrum itself rather than a scalar derived from it.

Answering that question is mainly a problem of measurement. When forecast quality is summarised
by an error aggregated over the whole spectrum, an error confined to a narrow peak (one or two
frequency bins out of several dozen) is diluted by the correctly predicted shoulders and tail. A forecast that merges distinct swell and wind-sea peaks into a single
smooth hump can therefore score as well as one that resolves them, or better, and an aggregate
criterion will prefer it. The effect has been seen before. A tendency to underestimate the
spectral peak was reported for an early neural spectral estimator and corrected there with an
empirical multiplier applied after the fact \citep{namekar2006ann}, and over-smoothing has been
identified as characteristic of data-driven forecasting models in meteorology
\citep{benbouallegue2024rise} and have already been reported in wave forecasting too \citep{minuzzi2023lstm}. The wave-modelling community has long recognised the underlying problem for NWMs. \citet{hanson2009pacific} note that mean or integral spectral properties give only
a general measure of performance and can mask higher-order deficiencies, and for that reason they
verify model output within individual wind-sea and swell partitions, an approach more recently applied to
long-range spectral forecasts \citep{rogers2025espc}. A further complication arises when the models
being compared were trained under different objectives. A model optimised for a given error measure tends
to score well on that measure, so a comparison based on either model's training objective cannot show
which model better represents the spectrum as a whole.

To our knowledge, no study has yet forecast the wave spectrum autoregressively from a single
station's own spectral history while treating the fidelity of the forecast spectral shape, and not
an aggregate error, as the quantity to optimise and report. This paper makes four contributions.
(i) We formulate spectral forecasting as sequence-to-sequence regression on the multi-channel
spectrum measured at a buoy. The model is a transformer whose input embedding keeps the frequency
axis as a structured dimension instead of flattening it, and whose output is parameterised
logarithmically, so the recovered spectrum is non-negative by construction and needs no
constrained output activation. (ii) We factorise the forecast target into a unit-area spectral
shape and a scalar magnitude, forecast the two separately and recombine them at prediction time.
This separates the comparatively easy magnitude problem from the harder shape problem. An
analogous factorisation is established in short-term load forecasting
\citep{sevlian2018scaling, guo2026loadshape}; here the magnitude is supplied by the physical
relation between wave height and spectral energy. (iii) We train against a composite spectral
objective combining a distributional divergence \citep{kullback1951}, an optimal-transport
distance \citep{peyre2019} and a differentiable peak-height term. Each targets a kind of spectral
error to which a per-bin criterion is insensitive, and model selection uses a criterion that does
not reduce to any one of them. (iv) We evaluate with a peak-resolved diagnostic panel in which spectral
peaks are identified by published significance criteria and classified as wind sea or swell \citep{portilla2009}, and forecast quality is reported separately for each class, so errors that an aggregate score would
dilute stay visible. Each component of (iii) is admitted on evidence through an incremental protocol: a
candidate is kept only if it improves the evaluation panel without degrading any other part of it,
and only if its measured computational cost justifies the gain.

The remainder of this paper is organised as follows. Section~\ref{sec:methods} presents the
methodology: the problem formulation and the shape/magnitude decomposition of the forecast target
(Section~\ref{sec:problem}), the buoy record and its preparation (Section~\ref{sec:data}), the
model architecture (Section~\ref{sec:architecture}), the training procedure and hyperparameter
search (Sections~\ref{sec:training} and~\ref{sec:hpo}), the evaluation protocol and its
peak-resolved, partition-conditioned diagnostics (Section~\ref{sec:evaluation}), and the
incremental ablation strategy under which each component was accepted or rejected
(Section~\ref{sec:ablation}). Section~\ref{sec:results} reports the results, and
Section~\ref{sec:discussion} discusses their implications and the limitations of the study.
